% ===== 末章：模型检验 / 灵敏度分析 / 模型评价与改进 / 附录 =====
% 对应 blueprints/paper_blueprint.md 章节树末端 0X_sensitivity_conclusion。
% 内容（按需保留小节，不必四节全有）：
%   模型检验（独立复算/交叉验证/约束扫描）→ 灵敏度与稳健性 → 模型评价与改进 → 附录（证明/结果文件/程序/复现）。

\section{模型检验}

\subsection{交叉验证或独立复算}
[用不同实现/离散化/方法复核关键数值，报告偏差；结论与正文一致。]

\subsection{约束扫描}
[路径/几何/过程约束在最易违背处独立抽查，报告违例数；引擎自报满足 ≠ 真实满足。]

\section{灵敏度与稳健性分析}

\subsection{关键参数敏感性}
[对关键参数做 ±5%/±10%/±20% 扰动，报告目标变化；指出最敏感参数。]

\subsection{口径/实现参数敏感性}
[判据口径、采样网格、时间步长等对结论影响；结论对这些参数是否稳健。]

\section{模型评价与改进方向}

\textbf{优点}：[列出 2-4 条具体优点，不空泛。]

\textbf{不足与改进}：[列出 2-4 条具体不足（符号复杂度、假设局限、未改进处），每条给改进方向，不做空泛自评。]

% ================= 附录（页数不计入正文 20 页上限） =================
\newpage

\section*{附录 A：理论性质}
[各问提炼的“精确归约”型性质（蓝图 §3.4 定理清单）的证明/梗概或“已尝试、不适用”说明。]

\section*{附录 B：结果文件说明}
[交付文件（result1/2/3.xlsx）的列结构、口径说明。]

\section*{附录 C：程序说明}
[求解引擎依赖、关键实现、报告值口径、复现命令。]

\section*{附录 D：核心代码节选}
\begin{lstlisting}[language=Python]
# 关键算法/原子算子节选（保持注释精简）
\end{lstlisting}

\section*{附录 E：复现步骤说明}
[一次运行复现全部结果的环境与命令序列。]
